% chapters/ch13-functional-dependencies.tex
\chapter{Functional Dependencies as a Design Tool}

\begin{learningobjectives}
In this chapter you will learn to use functional dependencies to justify schema structure.
You will be able to write dependencies from domain meaning, detect when a relation mixes facts that should be separated, and explain normalization decisions as correctness decisions rather than as formalisms.
\end{learningobjectives}

\section{Why dependencies matter}
When people first learn normalization, it often feels like a set of abstract rules.
In practice, normalization is a consequence of a simple idea: some facts determine others.
If you place facts together in a relation where the key does not determine them, you create redundancy.
That redundancy later becomes anomalies.

Functional dependencies are a language for stating “what determines what” precisely.
They are not only a theoretical tool.
They are a way to make your design decisions auditable.

\begin{definitionbox}
A \term{functional dependency} \(X \rightarrow Y\) means:
if two rows agree on \(X\), they must also agree on \(Y\).
In plain language: knowing \(X\) uniquely determines \(Y\).
\end{definitionbox}

You do not need heavy notation to benefit from this.
You need the habit of writing a few dependencies and checking whether your schema respects them.

\section{How to write dependencies from rules}
A common mistake is to infer dependencies from a small dataset.
Dependencies should come from domain rules, policies, and definitions.
They are claims about reality as modelled by the system.

\subsection{Example 1: identity-based dependency}
If \texttt{member\_id} is the identity of a member, then:
\[
member\_id \rightarrow name,\; email
\]
This is a statement about identity, not about observed data.

\subsection{Example 2: catalog dependency}
If ISBN identifies a title in your domain, then:
\[
isbn \rightarrow title,\; author
\]
If ISBN is optional or unreliable, this dependency may not hold, and the key choice should change.

\section{The dependency test for mixed relations}
Suppose you have a relation with key \(K\) and attributes \(A, B, C\).
If the meaning implies that some attribute \(B\) is determined by something other than \(K\), you are probably mixing facts.

This is the heart of normalization, stated as a practical test.

\begin{principlebox}
\textbf{Dependency test}\\
In a well-structured relation, non-key attributes should be determined by the key.
If a non-key attribute is determined by something else, you likely need a new relation.
\end{principlebox}

This is not a rigid law.
It is a strong signal.
When you break it deliberately, you should document why.

\begin{figure}[htbp]
\centering
\includegraphics[width=0.75\textwidth]{fd-decomposition.pdf}
\caption{Dependencies guide decomposition into relations whose keys determine their attributes.}
\end{figure}

\section{Worked example: clinic scheduling}
Consider a naive relation:
\[
\begin{aligned}
\rel{AppointmentRecord}(&\underline{appointment\_id}, start\_time, doctor\_name, \\
&doctor\_specialty, patient\_name)
\end{aligned}
\]
If \texttt{appointment\_id} is the key, this relation claims:
\[
appointment\_id \rightarrow doctor\_name,\; doctor\_specialty,\; patient\_name
\]
But from meaning, we also have:
\[
\begin{aligned}
doctor\_id &\rightarrow doctor\_name,\; doctor\_specialty \\
patient\_id &\rightarrow patient\_name
\end{aligned}
\]
So doctor facts and patient facts are being repeated across many appointments.
That repetition is what creates update anomalies.

The correct structural move is to separate entities:
\[
\rel{Doctor}(\underline{doctor\_id}, name, specialty)
\]
\[
\rel{Patient}(\underline{patient\_id}, name)
\]
\[
\begin{aligned}
\rel{Appointment}(&\underline{appointment\_id}, start\_time, \\
&doctor\_id \rightarrow \rel{Doctor}, patient\_id \rightarrow \rel{Patient})
\end{aligned}
\]

Now the dependencies align with keys:
\[
\begin{aligned}
doctor\_id &\rightarrow name,\; specialty \\
patient\_id &\rightarrow name \\
appointment\_id &\rightarrow start\_time,\; doctor\_id,\; patient\_id
\end{aligned}
\]

\section{Worked example: address modelling}
Suppose you store a customer’s address directly in \rel{Customer}:
\[
\rel{Customer}(\underline{customer\_id}, name, street, city, postal\_code)
\]
If customers can have multiple addresses, this structure is wrong.
The dependency is:
\[
customer\_id \nrightarrow street,\; city,\; postal\_code
\]
because a customer can have more than one address.

The correction is to introduce an association:
\[
\begin{aligned}
\rel{Address}(&\underline{address\_id}, street, city, postal\_code)
\end{aligned}
\]
\[
\begin{aligned}
\rel{CustomerAddress}(&\underline{customer\_id} \rightarrow \rel{Customer}, \\
&\underline{address\_id} \rightarrow \rel{Address}, \\
&address\_type)
\end{aligned}
\]
Now address facts depend on \texttt{address\_id}, and the association captures the relationship.

\section{Transitive dependencies and why they create redundancy}
A transitive dependency happens when:
\[
K \rightarrow A
\quad \text{and} \quad
A \rightarrow B
\]
so \(K \rightarrow B\) indirectly.

In design terms, this usually means you stored a code and its description in the same relation.
For example, storing \texttt{specialty\_name} inside \rel{Doctor} when \texttt{specialty\_id} exists and determines the name.

The redundancy is subtle at first, then painful later:
multiple spellings, inconsistent naming, and unclear governance.

The structural fix is to create a separate relation:
\[
\rel{Specialty}(\underline{specialty\_id}, name)
\]
and reference it from \rel{Doctor}.

\section{Dependencies and many-to-many relationships}
Dependencies also explain why many-to-many relationships require association relations.

If a doctor can have many specialties and a specialty can belong to many doctors, then neither of these holds:
\[
doctor\_id \rightarrow specialty\_id
\quad \text{nor} \quad
specialty\_id \rightarrow doctor\_id
\]
So the relationship itself must be represented as its own structure:
\[
\rel{DoctorSpecialty}(\underline{doctor\_id} \rightarrow \rel{Doctor}, \underline{specialty\_id} \rightarrow \rel{Specialty})
\]
This is not an implementation trick.
It is the correct representation of meaning.

\section{Dependencies as review language}
Dependencies are not just for normalization.
They give reviewers a precise language for critique.

When you say “this relation mixes facts,” you can support it with a dependency:
\[
A \rightarrow B \quad \text{but} \quad A \text{ is not a key}
\]
This turns a vague concern into a concrete design signal.

\begin{principlebox}
\textbf{Dependency review checklist}\\
For a relation \(\rel{R}(\underline{K}, A_1, \dots, A_n)\):
\begin{itemize}[leftmargin=*,nosep]
  \item Write 2--5 dependencies that come from meaning.
  \item Check that non-key attributes are determined by \(K\).
  \item If you find \(A \rightarrow B\) where \(A\) is not a key, ask whether you mixed facts.
\end{itemize}
\end{principlebox}

\begin{chaptersummary}
Functional dependencies express what determines what.
They come from domain meaning, not from sample data.
They provide a practical test for mixed relations: non-key attributes should be determined by the key.
When a non-key attribute is determined by something else, you are likely storing redundant facts and inviting anomalies.
\end{chaptersummary}

\begin{exercisebox}
Answer in full sentences and then write dependencies.
\begin{enumerate}[label=\arabic*.]
  \item Choose a domain and write a relation that you suspect is mixing facts.
  \item Propose a primary key and write the implied dependency \(K \rightarrow \text{(all other attributes)}\).
  \item Write two additional dependencies that come from meaning and show why redundancy exists.
  \item Restructure the schema into at least two relations and rewrite the dependencies so they align with keys.
  \item Give one example of a transitive dependency and show how you would eliminate it structurally.
\end{enumerate}
\end{exercisebox}
